% !PRINTED none

\begin{center}
\gesperrt{Sixth Section.}

\gesperrt{Tschirnhausen Transformation.}

\rule{4cm}{0.4pt}

\S.\,68.

\gesperrt{The Hermite Form of the Tschirnhausen Transformation.}
\end{center}

We have already become acquainted with the basic idea of the
Tschirnhausen transformation in the fourth section.

The task was this: to transform an algebraic equation of degree $n$
\weq{1}{f(x) = a_0 x^n + a_1 x^{n-1} + \cdots + a_{n-1} x + a_n = 0}
by a substitution of degree $(n-1)$
\weq{2}{y = \alpha_0 + \alpha_1 x + \alpha_2 x^2 + \cdots + \alpha_{n-1} x^{n-1}}
so as to gain, in the arbitrary coefficients $\alpha$, means of
simplifying the equation.

Hermite, by giving the substitution (2) a particular form, greatly
simplified this task and brought it into connection with the theory of
invariants\footnote{\gesperrt{Hermite}, Sur quelques th\'eor\`emes d'alg\`ebre et la r\'esolution de
l'\'equation du quatri\`eme degr\'e, published separately from the Comptes rendus of the Paris Academy. Paris 1859.}.

In \S.\,4 we became acquainted with a series of functions,
$f_0, f_1, f_2 \ldots f_{n-1}$, by which the powers of $x$ up to the
$(n-1)$st can be expressed rationally, and it is these functions
that Hermite uses for the representation of the substitution (2).

Here we denote by $x$ a root of the equation (1),
while $t$ is to be understood as an indeterminate variable.
Then by \S.\,4
%CONT>>
